% Options for packages loaded elsewhere
\PassOptionsToPackage{unicode}{hyperref}
\PassOptionsToPackage{hyphens}{url}
%
\documentclass[
]{article}
\usepackage{amsmath,amssymb}
\usepackage{iftex}
\ifPDFTeX
  \usepackage[T1]{fontenc}
  \usepackage[utf8]{inputenc}
  \usepackage{textcomp} % provide euro and other symbols
\else % if luatex or xetex
  \usepackage{unicode-math} % this also loads fontspec
  \defaultfontfeatures{Scale=MatchLowercase}
  \defaultfontfeatures[\rmfamily]{Ligatures=TeX,Scale=1}
\fi
\usepackage{lmodern}
\ifPDFTeX\else
  % xetex/luatex font selection
\fi
% Use upquote if available, for straight quotes in verbatim environments
\IfFileExists{upquote.sty}{\usepackage{upquote}}{}
\IfFileExists{microtype.sty}{% use microtype if available
  \usepackage[]{microtype}
  \UseMicrotypeSet[protrusion]{basicmath} % disable protrusion for tt fonts
}{}
\makeatletter
\@ifundefined{KOMAClassName}{% if non-KOMA class
  \IfFileExists{parskip.sty}{%
    \usepackage{parskip}
  }{% else
    \setlength{\parindent}{0pt}
    \setlength{\parskip}{6pt plus 2pt minus 1pt}}
}{% if KOMA class
  \KOMAoptions{parskip=half}}
\makeatother
\usepackage{xcolor}
\usepackage[margin=1in]{geometry}
\usepackage{color}
\usepackage{fancyvrb}
\newcommand{\VerbBar}{|}
\newcommand{\VERB}{\Verb[commandchars=\\\{\}]}
\DefineVerbatimEnvironment{Highlighting}{Verbatim}{commandchars=\\\{\}}
% Add ',fontsize=\small' for more characters per line
\usepackage{framed}
\definecolor{shadecolor}{RGB}{248,248,248}
\newenvironment{Shaded}{\begin{snugshade}}{\end{snugshade}}
\newcommand{\AlertTok}[1]{\textcolor[rgb]{0.94,0.16,0.16}{#1}}
\newcommand{\AnnotationTok}[1]{\textcolor[rgb]{0.56,0.35,0.01}{\textbf{\textit{#1}}}}
\newcommand{\AttributeTok}[1]{\textcolor[rgb]{0.13,0.29,0.53}{#1}}
\newcommand{\BaseNTok}[1]{\textcolor[rgb]{0.00,0.00,0.81}{#1}}
\newcommand{\BuiltInTok}[1]{#1}
\newcommand{\CharTok}[1]{\textcolor[rgb]{0.31,0.60,0.02}{#1}}
\newcommand{\CommentTok}[1]{\textcolor[rgb]{0.56,0.35,0.01}{\textit{#1}}}
\newcommand{\CommentVarTok}[1]{\textcolor[rgb]{0.56,0.35,0.01}{\textbf{\textit{#1}}}}
\newcommand{\ConstantTok}[1]{\textcolor[rgb]{0.56,0.35,0.01}{#1}}
\newcommand{\ControlFlowTok}[1]{\textcolor[rgb]{0.13,0.29,0.53}{\textbf{#1}}}
\newcommand{\DataTypeTok}[1]{\textcolor[rgb]{0.13,0.29,0.53}{#1}}
\newcommand{\DecValTok}[1]{\textcolor[rgb]{0.00,0.00,0.81}{#1}}
\newcommand{\DocumentationTok}[1]{\textcolor[rgb]{0.56,0.35,0.01}{\textbf{\textit{#1}}}}
\newcommand{\ErrorTok}[1]{\textcolor[rgb]{0.64,0.00,0.00}{\textbf{#1}}}
\newcommand{\ExtensionTok}[1]{#1}
\newcommand{\FloatTok}[1]{\textcolor[rgb]{0.00,0.00,0.81}{#1}}
\newcommand{\FunctionTok}[1]{\textcolor[rgb]{0.13,0.29,0.53}{\textbf{#1}}}
\newcommand{\ImportTok}[1]{#1}
\newcommand{\InformationTok}[1]{\textcolor[rgb]{0.56,0.35,0.01}{\textbf{\textit{#1}}}}
\newcommand{\KeywordTok}[1]{\textcolor[rgb]{0.13,0.29,0.53}{\textbf{#1}}}
\newcommand{\NormalTok}[1]{#1}
\newcommand{\OperatorTok}[1]{\textcolor[rgb]{0.81,0.36,0.00}{\textbf{#1}}}
\newcommand{\OtherTok}[1]{\textcolor[rgb]{0.56,0.35,0.01}{#1}}
\newcommand{\PreprocessorTok}[1]{\textcolor[rgb]{0.56,0.35,0.01}{\textit{#1}}}
\newcommand{\RegionMarkerTok}[1]{#1}
\newcommand{\SpecialCharTok}[1]{\textcolor[rgb]{0.81,0.36,0.00}{\textbf{#1}}}
\newcommand{\SpecialStringTok}[1]{\textcolor[rgb]{0.31,0.60,0.02}{#1}}
\newcommand{\StringTok}[1]{\textcolor[rgb]{0.31,0.60,0.02}{#1}}
\newcommand{\VariableTok}[1]{\textcolor[rgb]{0.00,0.00,0.00}{#1}}
\newcommand{\VerbatimStringTok}[1]{\textcolor[rgb]{0.31,0.60,0.02}{#1}}
\newcommand{\WarningTok}[1]{\textcolor[rgb]{0.56,0.35,0.01}{\textbf{\textit{#1}}}}
\usepackage{graphicx}
\makeatletter
\def\maxwidth{\ifdim\Gin@nat@width>\linewidth\linewidth\else\Gin@nat@width\fi}
\def\maxheight{\ifdim\Gin@nat@height>\textheight\textheight\else\Gin@nat@height\fi}
\makeatother
% Scale images if necessary, so that they will not overflow the page
% margins by default, and it is still possible to overwrite the defaults
% using explicit options in \includegraphics[width, height, ...]{}
\setkeys{Gin}{width=\maxwidth,height=\maxheight,keepaspectratio}
% Set default figure placement to htbp
\makeatletter
\def\fps@figure{htbp}
\makeatother
\setlength{\emergencystretch}{3em} % prevent overfull lines
\providecommand{\tightlist}{%
  \setlength{\itemsep}{0pt}\setlength{\parskip}{0pt}}
\setcounter{secnumdepth}{-\maxdimen} % remove section numbering
\ifLuaTeX
  \usepackage{selnolig}  % disable illegal ligatures
\fi
\usepackage{bookmark}
\IfFileExists{xurl.sty}{\usepackage{xurl}}{} % add URL line breaks if available
\urlstyle{same}
\hypersetup{
  pdftitle={STAT 7240: Homework 2},
  pdfauthor={{[}Your Name{]}},
  hidelinks,
  pdfcreator={LaTeX via pandoc}}

\title{STAT 7240: Homework 2}
\author{{[}Your Name{]}}
\date{September 25, 2024}

\begin{document}
\maketitle

\subsection{Load Required Libraries}\label{load-required-libraries}

\begin{Shaded}
\begin{Highlighting}[]
\FunctionTok{library}\NormalTok{(tidyverse)}
\FunctionTok{library}\NormalTok{(rpart)}
\FunctionTok{library}\NormalTok{(rpart.plot)}
\FunctionTok{library}\NormalTok{(ROCR)}
\FunctionTok{library}\NormalTok{(caret)}
\FunctionTok{library}\NormalTok{(ranger)}
\end{Highlighting}
\end{Shaded}

\subsection{Problem Set}\label{problem-set}

\subsubsection{Question 1 (2 points)}\label{question-1-2-points}

Read in the heart\_train.csv data as a data-frame. Conduct a quick
summary inspection of your data.

\begin{Shaded}
\begin{Highlighting}[]
\NormalTok{heart\_train }\OtherTok{\textless{}{-}} \FunctionTok{read.csv}\NormalTok{(}\StringTok{"heart\_train.csv"}\NormalTok{)}

\CommentTok{\# Summary inspection}
\FunctionTok{str}\NormalTok{(heart\_train)}
\end{Highlighting}
\end{Shaded}

\begin{verbatim}
## 'data.frame':    243 obs. of  14 variables:
##  $ age     : int  41 56 57 57 56 44 52 57 54 48 ...
##  $ sex     : int  0 1 0 1 0 1 1 1 1 0 ...
##  $ cp      : int  1 1 0 0 1 1 2 2 0 2 ...
##  $ trestbps: int  130 120 120 140 140 120 172 150 140 130 ...
##  $ chol    : int  204 236 354 192 294 263 199 168 239 275 ...
##  $ fbs     : int  0 0 0 0 0 0 1 0 0 0 ...
##  $ restecg : int  0 1 1 1 0 1 1 1 1 1 ...
##  $ thalach : int  172 178 163 148 153 173 162 174 160 139 ...
##  $ exang   : int  0 0 1 0 0 0 0 0 0 0 ...
##  $ oldpeak : num  1.4 0.8 0.6 0.4 1.3 0 0.5 1.6 1.2 0.2 ...
##  $ slope   : int  2 2 2 1 1 2 2 2 2 2 ...
##  $ ca      : int  0 0 0 0 0 0 0 0 0 0 ...
##  $ thal    : int  2 2 2 1 2 3 3 2 2 2 ...
##  $ target  : int  1 1 1 1 1 1 1 1 1 1 ...
\end{verbatim}

\begin{Shaded}
\begin{Highlighting}[]
\FunctionTok{summary}\NormalTok{(heart\_train)}
\end{Highlighting}
\end{Shaded}

\begin{verbatim}
##       age             sex               cp            trestbps    
##  Min.   :29.00   Min.   :0.0000   Min.   :0.0000   Min.   : 94.0  
##  1st Qu.:48.50   1st Qu.:0.0000   1st Qu.:0.0000   1st Qu.:120.0  
##  Median :56.00   Median :1.0000   Median :1.0000   Median :130.0  
##  Mean   :54.76   Mean   :0.6708   Mean   :0.9383   Mean   :131.7  
##  3rd Qu.:61.00   3rd Qu.:1.0000   3rd Qu.:2.0000   3rd Qu.:140.0  
##  Max.   :77.00   Max.   :1.0000   Max.   :3.0000   Max.   :192.0  
##       chol            fbs            restecg          thalach     
##  Min.   :131.0   Min.   :0.0000   Min.   :0.0000   Min.   : 71.0  
##  1st Qu.:211.5   1st Qu.:0.0000   1st Qu.:0.0000   1st Qu.:137.5  
##  Median :240.0   Median :0.0000   Median :1.0000   Median :154.0  
##  Mean   :246.8   Mean   :0.1523   Mean   :0.5185   Mean   :150.1  
##  3rd Qu.:273.5   3rd Qu.:0.0000   3rd Qu.:1.0000   3rd Qu.:166.0  
##  Max.   :564.0   Max.   :1.0000   Max.   :2.0000   Max.   :202.0  
##      exang           oldpeak          slope             ca        
##  Min.   :0.0000   Min.   :0.000   Min.   :0.000   Min.   :0.0000  
##  1st Qu.:0.0000   1st Qu.:0.000   1st Qu.:1.000   1st Qu.:0.0000  
##  Median :0.0000   Median :0.600   Median :1.000   Median :0.0000  
##  Mean   :0.2963   Mean   :1.002   Mean   :1.407   Mean   :0.7243  
##  3rd Qu.:1.0000   3rd Qu.:1.600   3rd Qu.:2.000   3rd Qu.:1.0000  
##  Max.   :1.0000   Max.   :6.200   Max.   :2.000   Max.   :4.0000  
##       thal           target      
##  Min.   :0.000   Min.   :0.0000  
##  1st Qu.:2.000   1st Qu.:0.0000  
##  Median :2.000   Median :1.0000  
##  Mean   :2.329   Mean   :0.5309  
##  3rd Qu.:3.000   3rd Qu.:1.0000  
##  Max.   :3.000   Max.   :1.0000
\end{verbatim}

\begin{Shaded}
\begin{Highlighting}[]
\FunctionTok{sum}\NormalTok{(}\FunctionTok{is.na}\NormalTok{(heart\_train))}
\end{Highlighting}
\end{Shaded}

\begin{verbatim}
## [1] 0
\end{verbatim}

\subsubsection{Question 2 (2 points)}\label{question-2-2-points}

Analyze the prevalence of heart disease in the dataset.

\begin{Shaded}
\begin{Highlighting}[]
\NormalTok{heart\_disease\_count }\OtherTok{\textless{}{-}} \FunctionTok{sum}\NormalTok{(heart\_train}\SpecialCharTok{$}\NormalTok{target }\SpecialCharTok{==} \DecValTok{1}\NormalTok{)}
\NormalTok{total\_count }\OtherTok{\textless{}{-}} \FunctionTok{nrow}\NormalTok{(heart\_train)}
\NormalTok{heart\_disease\_percentage }\OtherTok{\textless{}{-}}\NormalTok{ (heart\_disease\_count }\SpecialCharTok{/}\NormalTok{ total\_count) }\SpecialCharTok{*} \DecValTok{100}

\FunctionTok{print}\NormalTok{(}\FunctionTok{paste}\NormalTok{(}\StringTok{"Number of people with heart disease:"}\NormalTok{, heart\_disease\_count))}
\end{Highlighting}
\end{Shaded}

\begin{verbatim}
## [1] "Number of people with heart disease: 129"
\end{verbatim}

\begin{Shaded}
\begin{Highlighting}[]
\FunctionTok{print}\NormalTok{(}\FunctionTok{paste}\NormalTok{(}\StringTok{"Percentage of people with heart disease:"}\NormalTok{, }\FunctionTok{round}\NormalTok{(heart\_disease\_percentage, }\DecValTok{2}\NormalTok{), }\StringTok{"\%"}\NormalTok{))}
\end{Highlighting}
\end{Shaded}

\begin{verbatim}
## [1] "Percentage of people with heart disease: 53.09 %"
\end{verbatim}

\subsubsection{Question 3 (2 points)}\label{question-3-2-points}

Convert the target attribute to a factor.

\begin{Shaded}
\begin{Highlighting}[]
\NormalTok{heart\_train}\SpecialCharTok{$}\NormalTok{target }\OtherTok{\textless{}{-}} \FunctionTok{factor}\NormalTok{(heart\_train}\SpecialCharTok{$}\NormalTok{target, }\AttributeTok{levels =} \FunctionTok{c}\NormalTok{(}\DecValTok{1}\NormalTok{, }\DecValTok{0}\NormalTok{), }\AttributeTok{labels =} \FunctionTok{c}\NormalTok{(}\StringTok{"Yes"}\NormalTok{, }\StringTok{"No"}\NormalTok{))}
\end{Highlighting}
\end{Shaded}

\subsubsection{Question 4 (8 points)}\label{question-4-8-points}

Build and prune an rpart decision tree.

\begin{Shaded}
\begin{Highlighting}[]
\FunctionTok{set.seed}\NormalTok{(}\DecValTok{2}\NormalTok{)}

\CommentTok{\# Build the tree}
\NormalTok{tree\_model }\OtherTok{\textless{}{-}} \FunctionTok{rpart}\NormalTok{(target }\SpecialCharTok{\textasciitilde{}}\NormalTok{ ., }\AttributeTok{data =}\NormalTok{ heart\_train, }\AttributeTok{method =} \StringTok{"class"}\NormalTok{)}

\CommentTok{\# Find the optimal cp}
\NormalTok{optimal\_cp }\OtherTok{\textless{}{-}}\NormalTok{ tree\_model}\SpecialCharTok{$}\NormalTok{cptable[}\FunctionTok{which.min}\NormalTok{(tree\_model}\SpecialCharTok{$}\NormalTok{cptable[,}\StringTok{"xerror"}\NormalTok{]),}\StringTok{"CP"}\NormalTok{]}

\CommentTok{\# Prune the tree}
\NormalTok{pruned\_tree }\OtherTok{\textless{}{-}} \FunctionTok{prune}\NormalTok{(tree\_model, }\AttributeTok{cp =}\NormalTok{ optimal\_cp)}

\CommentTok{\# Plot the pruned tree}
\FunctionTok{rpart.plot}\NormalTok{(pruned\_tree, }\AttributeTok{extra =} \DecValTok{1}\NormalTok{, }\AttributeTok{under =} \ConstantTok{TRUE}\NormalTok{, }\AttributeTok{cex =} \FloatTok{0.8}\NormalTok{)}
\end{Highlighting}
\end{Shaded}

\includegraphics{Heartdiseaseanalysis_files/figure-latex/decision-tree-1.pdf}

\subsubsection{Question 5 (2 points)}\label{question-5-2-points}

considering the parameters there would be no changes required in the
controls.

\subsubsection{Question 6 (2 points)}\label{question-6-2-points}

from what they could be they could be important factors to determine
heart diseases based on age or cholestrol levels which then predicts the
risk of heart diseases

\subsubsection{Question 7 (10 points)}\label{question-7-10-points}

Predict on the test set and create a confusion matrix.

\begin{Shaded}
\begin{Highlighting}[]
\CommentTok{\# Read and prepare test data}
\NormalTok{heart\_test }\OtherTok{\textless{}{-}} \FunctionTok{read.csv}\NormalTok{(}\StringTok{"heart\_test.csv"}\NormalTok{)}
\NormalTok{heart\_test}\SpecialCharTok{$}\NormalTok{target }\OtherTok{\textless{}{-}} \FunctionTok{factor}\NormalTok{(heart\_test}\SpecialCharTok{$}\NormalTok{target, }\AttributeTok{levels =} \FunctionTok{c}\NormalTok{(}\DecValTok{1}\NormalTok{, }\DecValTok{0}\NormalTok{), }\AttributeTok{labels =} \FunctionTok{c}\NormalTok{(}\StringTok{"Yes"}\NormalTok{, }\StringTok{"No"}\NormalTok{))}

\CommentTok{\# Predict on test data}
\NormalTok{predictions }\OtherTok{\textless{}{-}} \FunctionTok{predict}\NormalTok{(pruned\_tree, heart\_test, }\AttributeTok{type =} \StringTok{"class"}\NormalTok{)}

\CommentTok{\# Create confusion matrix}
\NormalTok{conf\_matrix }\OtherTok{\textless{}{-}} \FunctionTok{table}\NormalTok{(}\AttributeTok{Predicted =}\NormalTok{ predictions, }\AttributeTok{Actual =}\NormalTok{ heart\_test}\SpecialCharTok{$}\NormalTok{target)}
\FunctionTok{print}\NormalTok{(conf\_matrix)}
\end{Highlighting}
\end{Shaded}

\begin{verbatim}
##          Actual
## Predicted Yes No
##       Yes  28  5
##       No    8 19
\end{verbatim}

\begin{Shaded}
\begin{Highlighting}[]
\CommentTok{\# Calculate accuracy}
\NormalTok{accuracy }\OtherTok{\textless{}{-}} \FunctionTok{sum}\NormalTok{(}\FunctionTok{diag}\NormalTok{(conf\_matrix)) }\SpecialCharTok{/} \FunctionTok{sum}\NormalTok{(conf\_matrix)}
\FunctionTok{print}\NormalTok{(}\FunctionTok{paste}\NormalTok{(}\StringTok{"Accuracy:"}\NormalTok{, }\FunctionTok{round}\NormalTok{(accuracy, }\DecValTok{4}\NormalTok{)))}
\end{Highlighting}
\end{Shaded}

\begin{verbatim}
## [1] "Accuracy: 0.7833"
\end{verbatim}

\subsubsection{Question 8 (2 points)}\label{question-8-2-points}

Analyze False Positives vs False Negatives.

\begin{Shaded}
\begin{Highlighting}[]
\NormalTok{false\_positives }\OtherTok{\textless{}{-}}\NormalTok{ conf\_matrix[}\StringTok{"Yes"}\NormalTok{, }\StringTok{"No"}\NormalTok{]}
\NormalTok{false\_negatives }\OtherTok{\textless{}{-}}\NormalTok{ conf\_matrix[}\StringTok{"No"}\NormalTok{, }\StringTok{"Yes"}\NormalTok{]}

\FunctionTok{print}\NormalTok{(}\FunctionTok{paste}\NormalTok{(}\StringTok{"False Positives:"}\NormalTok{, false\_positives))}
\end{Highlighting}
\end{Shaded}

\begin{verbatim}
## [1] "False Positives: 5"
\end{verbatim}

\begin{Shaded}
\begin{Highlighting}[]
\FunctionTok{print}\NormalTok{(}\FunctionTok{paste}\NormalTok{(}\StringTok{"False Negatives:"}\NormalTok{, false\_negatives))}
\end{Highlighting}
\end{Shaded}

\begin{verbatim}
## [1] "False Negatives: 8"
\end{verbatim}

\subsubsection{Question 9 (8 points)}\label{question-9-8-points}

Create an ROC curve.

\begin{Shaded}
\begin{Highlighting}[]
\CommentTok{\# Get probabilities}
\NormalTok{probs }\OtherTok{\textless{}{-}} \FunctionTok{predict}\NormalTok{(pruned\_tree, heart\_test, }\AttributeTok{type =} \StringTok{"prob"}\NormalTok{)[,}\StringTok{"Yes"}\NormalTok{]}

\CommentTok{\# Create prediction object}
\NormalTok{pred }\OtherTok{\textless{}{-}} \FunctionTok{prediction}\NormalTok{(probs, heart\_test}\SpecialCharTok{$}\NormalTok{target)}

\CommentTok{\# Create performance object}
\NormalTok{perf }\OtherTok{\textless{}{-}} \FunctionTok{performance}\NormalTok{(pred, }\StringTok{"tpr"}\NormalTok{, }\StringTok{"fpr"}\NormalTok{)}

\CommentTok{\# Create data frame for ggplot}
\NormalTok{roc\_data }\OtherTok{\textless{}{-}} \FunctionTok{data.frame}\NormalTok{(}\AttributeTok{FPR =}\NormalTok{ perf}\SpecialCharTok{@}\NormalTok{x.values[[}\DecValTok{1}\NormalTok{]], }\AttributeTok{TPR =}\NormalTok{ perf}\SpecialCharTok{@}\NormalTok{y.values[[}\DecValTok{1}\NormalTok{]])}

\CommentTok{\# Plot ROC curve}
\FunctionTok{ggplot}\NormalTok{(roc\_data, }\FunctionTok{aes}\NormalTok{(}\AttributeTok{x =}\NormalTok{ FPR, }\AttributeTok{y =}\NormalTok{ TPR)) }\SpecialCharTok{+}
  \FunctionTok{geom\_line}\NormalTok{() }\SpecialCharTok{+}
  \FunctionTok{geom\_abline}\NormalTok{(}\AttributeTok{intercept =} \DecValTok{0}\NormalTok{, }\AttributeTok{slope =} \DecValTok{1}\NormalTok{, }\AttributeTok{linetype =} \StringTok{"dashed"}\NormalTok{) }\SpecialCharTok{+}
  \FunctionTok{labs}\NormalTok{(}\AttributeTok{title =} \StringTok{"ROC Curve for Heart Disease Prediction"}\NormalTok{,}
       \AttributeTok{x =} \StringTok{"False Positive Rate"}\NormalTok{, }
       \AttributeTok{y =} \StringTok{"True Positive Rate"}\NormalTok{) }\SpecialCharTok{+}
  \FunctionTok{theme\_minimal}\NormalTok{()}
\end{Highlighting}
\end{Shaded}

\includegraphics{Heartdiseaseanalysis_files/figure-latex/roc-curve-1.pdf}

\subsubsection{Question 10 (2 points)}\label{question-10-2-points}

Find the AUC of your ROC curve.

\begin{Shaded}
\begin{Highlighting}[]
\NormalTok{auc }\OtherTok{\textless{}{-}} \FunctionTok{performance}\NormalTok{(pred, }\StringTok{"auc"}\NormalTok{)}\SpecialCharTok{@}\NormalTok{y.values[[}\DecValTok{1}\NormalTok{]]}
\FunctionTok{print}\NormalTok{(}\FunctionTok{paste}\NormalTok{(}\StringTok{"AUC:"}\NormalTok{, }\FunctionTok{round}\NormalTok{(auc, }\DecValTok{4}\NormalTok{)))}
\end{Highlighting}
\end{Shaded}

\begin{verbatim}
## [1] "AUC: 0.8194"
\end{verbatim}

\subsubsection{Question 11 (5 points)}\label{question-11-5-points}

Plot Accuracy vs Cutoff.

\begin{Shaded}
\begin{Highlighting}[]
\NormalTok{acc }\OtherTok{\textless{}{-}} \FunctionTok{performance}\NormalTok{(pred, }\StringTok{"acc"}\NormalTok{)}
\NormalTok{acc\_data }\OtherTok{\textless{}{-}} \FunctionTok{data.frame}\NormalTok{(}\AttributeTok{Cutoff =}\NormalTok{ acc}\SpecialCharTok{@}\NormalTok{x.values[[}\DecValTok{1}\NormalTok{]], }\AttributeTok{Accuracy =}\NormalTok{ acc}\SpecialCharTok{@}\NormalTok{y.values[[}\DecValTok{1}\NormalTok{]])}

\FunctionTok{ggplot}\NormalTok{(acc\_data, }\FunctionTok{aes}\NormalTok{(}\AttributeTok{x =}\NormalTok{ Cutoff, }\AttributeTok{y =}\NormalTok{ Accuracy)) }\SpecialCharTok{+}
  \FunctionTok{geom\_line}\NormalTok{() }\SpecialCharTok{+}
  \FunctionTok{labs}\NormalTok{(}\AttributeTok{title =} \StringTok{"Accuracy vs Cutoff"}\NormalTok{,}
       \AttributeTok{x =} \StringTok{"Cutoff"}\NormalTok{, }
       \AttributeTok{y =} \StringTok{"Accuracy"}\NormalTok{) }\SpecialCharTok{+}
  \FunctionTok{theme\_minimal}\NormalTok{()}
\end{Highlighting}
\end{Shaded}

\includegraphics{Heartdiseaseanalysis_files/figure-latex/accuracy-vs-cutoff-1.pdf}

\subsubsection{Question 12 (15 points)}\label{question-12-15-points}

Tune a random forest model using the caret package.

\begin{Shaded}
\begin{Highlighting}[]
\FunctionTok{set.seed}\NormalTok{(}\DecValTok{2}\NormalTok{)}

\CommentTok{\# Define hyperparameters}
\NormalTok{hyper\_grid }\OtherTok{\textless{}{-}} \FunctionTok{expand.grid}\NormalTok{(}
  \AttributeTok{mtry =} \FunctionTok{c}\NormalTok{(}\DecValTok{2}\NormalTok{, }\DecValTok{4}\NormalTok{, }\DecValTok{6}\NormalTok{),}
  \AttributeTok{splitrule =} \StringTok{"gini"}\NormalTok{,}
  \AttributeTok{min.node.size =} \FunctionTok{c}\NormalTok{(}\DecValTok{1}\NormalTok{, }\DecValTok{3}\NormalTok{, }\DecValTok{5}\NormalTok{)}
\NormalTok{)}

\CommentTok{\# Define training control}
\NormalTok{ctrl }\OtherTok{\textless{}{-}} \FunctionTok{trainControl}\NormalTok{(}
  \AttributeTok{method =} \StringTok{"cv"}\NormalTok{,}
  \AttributeTok{number =} \DecValTok{10}\NormalTok{,}
  \AttributeTok{classProbs =} \ConstantTok{TRUE}\NormalTok{,}
  \AttributeTok{summaryFunction =}\NormalTok{ twoClassSummary}
\NormalTok{)}

\CommentTok{\# Train the model}
\NormalTok{rf\_model }\OtherTok{\textless{}{-}} \FunctionTok{train}\NormalTok{(}
\NormalTok{  target }\SpecialCharTok{\textasciitilde{}}\NormalTok{ .,}
  \AttributeTok{data =}\NormalTok{ heart\_train,}
  \AttributeTok{method =} \StringTok{"ranger"}\NormalTok{,}
  \AttributeTok{num.trees =} \DecValTok{500}\NormalTok{,}
  \AttributeTok{importance =} \StringTok{"impurity"}\NormalTok{,}
  \AttributeTok{tuneGrid =}\NormalTok{ hyper\_grid,}
  \AttributeTok{trControl =}\NormalTok{ ctrl,}
  \AttributeTok{metric =} \StringTok{"ROC"}
\NormalTok{)}

\CommentTok{\# Print best model}
\FunctionTok{print}\NormalTok{(rf\_model}\SpecialCharTok{$}\NormalTok{bestTune)}
\end{Highlighting}
\end{Shaded}

\begin{verbatim}
##   mtry splitrule min.node.size
## 2    2      gini             3
\end{verbatim}

\begin{Shaded}
\begin{Highlighting}[]
\CommentTok{\# Predict on test set}
\NormalTok{rf\_predictions }\OtherTok{\textless{}{-}} \FunctionTok{predict}\NormalTok{(rf\_model, heart\_test)}
\NormalTok{rf\_probs }\OtherTok{\textless{}{-}} \FunctionTok{predict}\NormalTok{(rf\_model, heart\_test, }\AttributeTok{type =} \StringTok{"prob"}\NormalTok{)}

\CommentTok{\# Create confusion matrix}
\NormalTok{rf\_conf\_matrix }\OtherTok{\textless{}{-}} \FunctionTok{confusionMatrix}\NormalTok{(rf\_predictions, heart\_test}\SpecialCharTok{$}\NormalTok{target)}
\FunctionTok{print}\NormalTok{(rf\_conf\_matrix)}
\end{Highlighting}
\end{Shaded}

\begin{verbatim}
## Confusion Matrix and Statistics
## 
##           Reference
## Prediction Yes No
##        Yes  30  4
##        No    6 20
##                                           
##                Accuracy : 0.8333          
##                  95% CI : (0.7148, 0.9171)
##     No Information Rate : 0.6             
##     P-Value [Acc > NIR] : 8.893e-05       
##                                           
##                   Kappa : 0.6575          
##                                           
##  Mcnemar's Test P-Value : 0.7518          
##                                           
##             Sensitivity : 0.8333          
##             Specificity : 0.8333          
##          Pos Pred Value : 0.8824          
##          Neg Pred Value : 0.7692          
##              Prevalence : 0.6000          
##          Detection Rate : 0.5000          
##    Detection Prevalence : 0.5667          
##       Balanced Accuracy : 0.8333          
##                                           
##        'Positive' Class : Yes             
## 
\end{verbatim}

\subsubsection{Bonus Question (2 points)}\label{bonus-question-2-points}

Compare the accuracy of your tuned Random Forest model to your pruned
Decision Tree on the test sample.

\begin{Shaded}
\begin{Highlighting}[]
\NormalTok{dt\_accuracy }\OtherTok{\textless{}{-}}\NormalTok{ accuracy  }\CommentTok{\# from Question 7}
\NormalTok{rf\_accuracy }\OtherTok{\textless{}{-}}\NormalTok{ rf\_conf\_matrix}\SpecialCharTok{$}\NormalTok{overall[}\StringTok{"Accuracy"}\NormalTok{]}

\FunctionTok{print}\NormalTok{(}\FunctionTok{paste}\NormalTok{(}\StringTok{"Decision Tree Accuracy:"}\NormalTok{, }\FunctionTok{round}\NormalTok{(dt\_accuracy, }\DecValTok{4}\NormalTok{)))}
\end{Highlighting}
\end{Shaded}

\begin{verbatim}
## [1] "Decision Tree Accuracy: 0.7833"
\end{verbatim}

\begin{Shaded}
\begin{Highlighting}[]
\FunctionTok{print}\NormalTok{(}\FunctionTok{paste}\NormalTok{(}\StringTok{"Random Forest Accuracy:"}\NormalTok{, }\FunctionTok{round}\NormalTok{(rf\_accuracy, }\DecValTok{4}\NormalTok{)))}
\end{Highlighting}
\end{Shaded}

\begin{verbatim}
## [1] "Random Forest Accuracy: 0.8333"
\end{verbatim}

\end{document}
